\chapter*{Введение}
\label{sec:afterwords}
\addcontentsline{toc}{chapter}{Введение}

Ролевые игры традиционно требуют участия мастера игры, который управляет миром, неигровыми персонажами и развитием сюжета. Развитие технологий больших языковых моделей открывает возможности для автоматизации ролевых игр, однако существующие коммерческие решения (AI Dungeon, AI Dungeon-Chi и др.) основаны на чисто генеративном подходе без структурированного управления состоянием игры, что приводит к проблемам согласованности повествования и отсутствию контроля действий игрока.

Новизна работы заключается в интеграции системы управления эмоциями персонажей на основе моральных схем, механизма директора для координации действий персонажей и системы контроля действий в единую архитектуру, обеспечивающую структурированное управление игровым процессом.

Оригинальная суть исследования состоит в создании системы, где виртуальный мастер и неигровые персонажи управляются искусственным интеллектом, при этом каждый неигровой персонаж обладает собственной моделью эмоций и намерений, которая обновляется в процессе игры, а виртуальный мастер координирует их действия для поддержания связности сюжета.

В первой главе проводится анализ существующих подходов к созданию систем виртуальных мастеров, рассматриваются коммерческие решения и их проблемы, выполняется сравнительный анализ различных систем.

Во второй главе описываются разработанные модели и алгоритмы: формальная модель системы ролевой игры, модель управления эмоциональным состоянием персонажей, алгоритмы работы основных компонентов и обобщенная архитектура системы.

В третьей главе описываются результаты проектирования: архитектура основных подсистем (управление игровым циклом, сценами, выдача директив, управление персонажами, эмоциональное моделирование), их интерфейсы и взаимодействие.

В четвертой главе описывается реализация системы на языке Python, выбор инструментальных средств, структура программного обеспечения, примеры работы и результаты тестирования.


%%% Local Variables:
%%% TeX-engine: xetex
%%% eval: (setq-local TeX-master (concat "../" (seq-find (-cut string-match ".*-3-pz\.tex$" <>) (directory-files ".."))))
%%% End:
